\documentclass[11pt,a4paper]{article}
\usepackage[utf8]{inputenc}
\usepackage[T1]{fontenc} % Necesario para guillemots y caracteres especiales
\usepackage[spanish]{babel}
\usepackage{tikz}
% Añadida shapes.geometric para el cilindro (database) y arrows.meta para flechas modernas
\usetikzlibrary{shapes.geometric,arrows.meta,positioning,fit,backgrounds,calc,babel}
\usepackage{geometry}
\geometry{a4paper, margin=2cm}
\usepackage{parskip}

\title{\textbf{Documentación de la Arquitectura}\\ \Large Práctica 3 - Broker de Mensajes}
\author{Imad Habib Jali (901421) \\ Jorge Bellido Lobera (903080)}
\date{\today}

\begin{document}

\maketitle

\section{Introducción}
El presente documento describe la arquitectura del \textit{Message-Oriented Middleware (MOM)} desarrollado. Se presentan tres vistas arquitectónicas fundamentales para comprender la estructura, interacción y despliegue del sistema.

\section{1. Vista de Módulos}
Ilustra la estructura estática del código y las dependencias entre módulos Python.

\begin{center}
\begin{tikzpicture}[
    module/.style={draw, rectangle, rounded corners, fill=blue!10, text centered, text width=3cm, minimum height=1.5cm, font=\sffamily\small},
    dependency/.style={-{Stealth}, dashed, thick}
]

% Nodos
\node[module] (client) {client.py \\ \footnotesize(Librería Cliente)};
\node[module, above left=1.5cm and 1cm of client] (productor) {productor.py \\ \footnotesize(App Productora)};
\node[module, above right=1.5cm and 1cm of client] (consumidor) {consumidor.py \\ \footnotesize(App Consumidora)};
\node[module, below=2cm of client] (server) {server.py \\ \footnotesize(Broker Core)};
\node[module, below=1.5cm of server] (storage) {storage.py \\ \footnotesize(Persistencia)};

% Conexiones
\draw[dependency] (productor) -- (client);
\draw[dependency] (consumidor) -- (client);
\draw[dependency] (server) -- (storage);

% Fronteras
\begin{scope}[on background layer]
    \node[draw, dashed, fill=gray!5, fit=(productor) (consumidor) (client), inner sep=0.5cm, label={[anchor=north west]north west:\textbf{Lado Cliente}}] {};
    \node[draw, dashed, fill=gray!5, fit=(server) (storage), inner sep=0.5cm, label={[anchor=north west]north west:\textbf{Lado Servidor}}] {};
\end{scope}

\end{tikzpicture}
\end{center}

\newpage
\section{2. Vista de Componentes y Conectores}
Describe el sistema en tiempo de ejecución, mostrando los procesos y el flujo de mensajes.

\begin{center}
\begin{tikzpicture}[
    component/.style={draw, rectangle, fill=green!10, text centered, text width=2.5cm, minimum height=1.5cm, font=\sffamily\scriptsize},
    interface/.style={draw, circle, fill=black, inner sep=1.5pt},
    connector/.style={-{Stealth}, thick, color=blue!70!black},
    database/.style={draw, cylinder, shape border rotate=90, aspect=0.25, fill=yellow!20, minimum height=1.2cm, minimum width=1.8cm, align=center, text centered, font=\sffamily\tiny}
]

% Procesos externos
\node[component] (prodComp) {Proceso Productor};
\node[component, below=3.5cm of prodComp] (consComp) {Proceso Consumidor};

% Broker (Contenedor más estrecho y distancia reducida)
\node[draw, rectangle, fill=red!5, right=2cm of prodComp, text width=4.5cm, minimum height=6.5cm, anchor=north west] (brokerContainer) at ($(prodComp.north east)+(2cm,0.5cm)$) {};
\node[anchor=north, font=\sffamily\scriptsize\bfseries] at (brokerContainer.north) {Broker (Servidor)};

\node[draw, fill=white, text width=3cm, align=center, font=\sffamily\tiny] (tcpHandler) at ($(brokerContainer.center)+(0,1.8cm)$) {Manejador TCP \\ (Fair Dispatch / ACKs)};
\node[draw, fill=white, text width=3cm, align=center, font=\sffamily\tiny] (queueMem) at (brokerContainer.center) {Colas en Memoria \\ (Diccionario RLock)};
\node[draw, fill=white, text width=3cm, align=center, font=\sffamily\tiny] (ttlThread) at ($(brokerContainer.center)-(0,1.8cm)$) {Hilo Cleanup \\ (TTL / ACK Timeouts)};

\draw[{Stealth}-{Stealth}, dashed] (tcpHandler) -- (queueMem);
\draw[-{Stealth}, dashed] (ttlThread) -- (queueMem);

% Almacenamiento (Más cerca del broker)
\node[database, right=0.8cm of queueMem] (db) {JSON File \\ storage};

% Puerto
\node[interface, left=0.3cm of tcpHandler, label=above:\tiny Port 5555] (port) {};

% Conectores
\draw[connector] (prodComp.east) -- node[above, sloped, font=\tiny] {Publish} (port);
\draw[connector] (port) -- node[below, sloped, font=\tiny] {Push / ACK} (consComp.east);
\draw[connector, dashed] (queueMem.east) -- (db.west);

\end{tikzpicture}
\end{center}

\section{3. Vista de Distribución e Instalación}
Esta vista detalla el despliegue del sistema mediante tres niveles de abstracción física.

\subsection{3.1. Vista de Asignación (Componentes a Artefactos)}
Mapeo de la lógica del sistema a los archivos ejecutables reales.

\begin{center}
\begin{tikzpicture}[
    comp/.style={draw, rectangle, fill=blue!5, text width=2.5cm, align=center, font=\sffamily\tiny},
    art/.style={draw, rectangle, fill=white, text width=2.2cm, align=center, font=\sffamily\tiny, double}
]
\node[comp] (c1) {Componente\\Broker Core};
\node[comp, right=2cm of c1] (c2) {Componente\\Persistencia};

\node[art, below=1cm of c1] (a1) {<<artifact>>\\server.py};
\node[art, below=1cm of c2] (a2) {<<artifact>>\\storage.py};

% Conexiones de asignación
\draw[dashed, ->] (c1) -- (a1);
\draw[dashed, ->] (c2) -- (a2);

% Dependencia entre componentes
\draw[->, thick] (c1) -- node[above, font=\tiny] {usa} (c2);
\end{tikzpicture}
\end{center}

\subsection{3.2. Vista de Despliegue (Artefactos a Nodos)}
Configuración de la red y ubicación de los artefactos en los nodos de ejecución.

\begin{center}
\begin{tikzpicture}[
    nodeOS/.style={draw, thick, fill=gray!10, rectangle, minimum width=3.5cm, minimum height=2.5cm, align=center, font=\sffamily\tiny},
    art/.style={draw, fill=white, rectangle, text width=2.5cm, align=center, font=\sffamily\tiny}
]
% Nodo Servidor (Efecto 3D)
\node[nodeOS] (nodeA) at (0,0) {\textbf{Nodo: Servidor}\\ \footnotesize(Ubuntu/Windows)};
\draw (nodeA.north west) -- ++(0.2,0.2) -- ($(nodeA.north east)+(0.2,0.2)$) -- (nodeA.north east);
\draw (nodeA.north east) -- ++(0.2,0.2) -- ($(nodeA.south east)+(0.2,0.2)$) -- (nodeA.south east);

\node[art, below=0.2cm of nodeA.center, anchor=north] (artS) {<<artifact>>\\server.py / storage.py};

% Nodos Clientes
\node[nodeOS] (nodeB) at (-5, -4) {\textbf{Nodo: Emisor}};
\node[art, below=0.5cm of nodeB.north] {<<artifact>>\\client.py};

\node[nodeOS] (nodeC) at (5, -4) {\textbf{Nodo: Receptor}};
\node[art, below=0.5cm of nodeC.north] {<<artifact>>\\client.py};

\draw[thick, <->] (nodeB) -- node[above, sloped, font=\tiny] {TCP:5555} (nodeA);
\draw[thick, <->] (nodeC) -- node[above, sloped, font=\tiny] {TCP:5555} (nodeA);
\end{tikzpicture}
\end{center}

\subsection{3.3. Vista de Instalación (Paquete de Aplicación)}
Estructura física completa de los componentes, recursos y datos necesarios para el funcionamiento del sistema en los entornos de destino.

\begin{center}
\begin{tikzpicture}[
    node distance=0.5cm,
    package/.style={draw, rectangle, fill=yellow!5, text width=8cm, align=left, font=\sffamily\tiny, thick},
    module/.style={font=\sffamily\tiny, inner sep=1pt},
    data/.style={font=\sffamily\tiny\itshape, color=gray!80!black}
]

% Paquete Servidor
\node[package] (serverPkg) {
    \textbf{Directorio del Servidor (/broker-srv/)} \\
    \vspace{0.1cm}
    \begin{tabular}{l l}
    \module $\bullet$ \texttt{server.py} & \module Núcleo del Message Broker (Gestión de colas y red) \\
    \module $\bullet$ \texttt{storage.py} & \module Gestor de Persistencia Atómica (I/O Disco) \\
    \data $\bullet$ \texttt{broker\_storage.json} & \data Base de datos persistente (Estado de colas y mensajes) \\
    \data $\bullet$ \texttt{broker.log} & \data Trazas de ejecución, auditoría y errores del sistema \\
    \end{tabular}
};

% Paquete Cliente
\node[package, below=1cm of serverPkg] (clientPkg) {
    \textbf{Directorio de Aplicación Cliente (/client-app/)} \\
    \vspace{0.1cm}
    \begin{tabular}{l l}
    \module $\bullet$ \texttt{client.py} & \module Librería Cliente (API de comunicación JSON sobre TCP) \\
    \module $\bullet$ \texttt{productor.py} & \module Lógica de aplicación para publicación de mensajes \\
    \module $\bullet$ \texttt{consumidor.py} & \module Lógica de aplicación para procesamiento de mensajes \\
    \end{tabular}
};

% Recursos Comunes
\node[package, below=1cm of clientPkg] (docPkg) {
    \textbf{Documentación y Pruebas} \\
    \vspace{0.1cm}
    \begin{tabular}{l l}
    \module $\bullet$ \texttt{test\_main.py} & \module Suite de Pruebas de Integración y Stress (6 casos) \\
    \module $\bullet$ \texttt{arquitectura.pdf} & \module Documento de diseño arquitectónico (Vistas 1-3) \\
    \end{tabular}
};

\draw[thick, gray] (serverPkg.west) -- ++(-0.4,0) |- (clientPkg.west);
\draw[thick, gray] (clientPkg.west) -- ++(-0,0) |- (docPkg.west);

\end{tikzpicture}
\end{center}

\end{document}
